% !TEX root = cap03_standalone.tex
% Capítulo 3. Estado del arte
% Reestructurado siguiendo la revisión de la Dra. Lárraga:
% catálogo descriptivo -> análisis crítico por problemas metodológicos.

\chapter{Estado del arte}
\label{ch:estado-del-arte}

\begin{resumen}
El modelado computacional del crecimiento urbano ha producido en treinta años
una literatura rica pero metodológicamente diferente. Este capítulo no busca
ordenar trabajos por autor sino que busca identificar los cuatro problemas
metodológicos que siguen abiertos en el campo, calibración por fuerza bruta,
validación sobreajustada a una sola ventana temporal, opacidad de la función de
transición en modelos de aprendizaje profundo y dependencia de software
propietario, no se promueve el software libre, y muestra cómo cada uno justifica una decisión concreta del diseño
del modelo WoE-AC propuesto. La discusión concluye con un mapa de la literatura
para ciudades mexicanas, donde ningún trabajo publicado satisface simultáneamente
los requerimientos R1--R4 de la Tabla~\ref{tab:requerimientos_modelo_1}.
\end{resumen}

% ============================================================
\section{Evolución del modelado urbano computacional}
\label{sec:evolucion}
% ============================================================

El modelado de la dinámica urbana mediante autómatas celulares empezó con
\textcite{White1997}, cuyo trabajo mostró que reglas locales basadas en distancia y
vecindad podrían reproducir patrones macroscópicos realistas en ciudades. El siguiente nivel
del enfoque fue introducido con SLEUTH en \parencite{Clarke1998,Silva2002}, que durante dos
décadas fue vanguardia y considerado como referencia, y que se aplicó en decenas de ciudades.

Al mismo tiempo, la metodología LUCC (Land Use and Cover Change) introdujo modelos
basados en mapas de el mejor estado posible de uso de suelo estadístico: sistemas como CLUE-S, Dinamica EGO,
CA-Markov y, más recientemente, FLUS, GeoSOS y MOLUSCE, brindaron
arquitecturas alternativas al autómata celular puro en las que la función de transición
emerge de regresiones logísticas o pesos bayesianos sobre datos históricos.

Una conjunto de académicos y profesionales de la planeación urbana posterior sustituyó la regresión por modelos de aprendizaje
automático e híbridos metaheurísticos. \textcite{almeida2008stochastic} utilizó
una red neuronal feedforward a un autómata celular estocástico para Piracicaba, Brasil; años
más tarde \textcite{Wang2021} combinó un autómata celular con un algoritmo
genético para optimizar la forma urbana de Wuhan (la busqueda fue ineficiente en términos de tiempo computacional), mientras
\textcite{Hagenauer2022} introdujo una red neuronal geográficamente ponderada
para capturar relaciones espaciales no lineales. Estos enfoques reportan
ganancias cuantitativas en temrinos de precisión, pero sustituyen parámetros legibles (o que pueden ser explicados) por activaciones
internas o cromosomas que no admiten una lectura geográfica directa, es decir, no son auditables en cuestión de lo que aprenden.

La decisión central del campo no es en torno a AC y aprendizaje automático,
sino entre dos directivas que a priori parecen contrapuestos: maximizar precisión
predictiva y mantener la trazabilidad de cada decisión espacial. Las secciones
siguientes muestran que esta tensión es resoluble si se cambia la pregunta
metodológica.

% ============================================================
\section{Problemas metodológicos persistentes}
\label{sec:problemas-metodologicos}
% ============================================================

Los problemas metodológicos mas importantes detectados en la literatura fueron 
equifinalidad en la calibración, sobreajuste en la validación, opacidad en la función de transición y dependencia de software no libre.

\subsection{Calibración: equifinalidad y costo computacional}

SLEUTH calibra cinco coeficientes enteros mediante búsqueda en cuadrícula sobre
el espacio discreto $[0,100]^4 \times [0,20]$, es decir, en el orden de $10^9$
combinaciones evaluadas con simulaciones Monte Carlo de 100 iteraciones cada
una. \textcite{Clarke1998} resolvió el problema con tres fases: coarse, fine,
final, que reducen el espacio de búsqueda iterativamente; aun así, una calibración completa
toma días en hardware contemporáneo \parencite{Jantz2003}. El problema
metodológico no es solo el costo: es la \textbf{equifinalidad}, esto es, la
existencia de múltiples combinaciones de coeficientes que producen patrones
indistinguibles a la métrica de calibración pero divergen en proyecciones de
largo plazo \parencite{Clarke2008}. Sin un criterio adicional, la elección
entre combinaciones equifinales termina dependiendo de la persona que realice la calibración.

El presente trabajo evita el problema redefiniéndolo: en lugar de calibrar
coeficientes que controlan reglas geométricas genéricas (diffusion, breed,
spread), calibra dos parámetros que controlan la combinación de evidencia
espacial empírica (umbral de transición y peso de vecindad). El espacio de
búsqueda es de dos dimensiones y la calibración termina en minutos. Más
importante: cada configuración candidata produce mapas de estado deseado cuyos pesos
por variable son directamente auditables o sabemos por qué tomó una decision, dicho de otra forma, eliminando la ambigüedad asociada a
la equifinalidad.

\subsection{Validación: sobreajuste temporal y métricas no comparables}

\textcite{pontius2008comparing} analizó trece simulaciones independientes de
nueve modelos de cambio de uso de suelo y propuso el Figure of Merit (FoM) como
métrica estándar comparable entre estudios. Su hallazgo es un tanto desmotivante: seis de
las trece aplicaciones obtienen FoM inferior al $15\%$ y solo una supera el
$50\%$. La dispersión o variabilidad no se explica por la complejidad del modelo sino por la
magnitud del cambio neto observado: el FoM crece monotónicamente con el
porcentaje de celdas que efectivamente cambian de estado. O dicho de una forma mas sencilla, el FoM es una métrica que penaliza la cantidad de celdas que no cambian de estado.

Esto tiene dos implicaciones que la literatura rara vez confronta. La primera,
\textbf{validar en una sola ventana temporal es estadísticamente insuficiente}:
si el FoM depende de la magnitud del cambio, entonces se puede concluir que una validación única no permite
distinguir entre un modelo robusto y uno sobreajustado a una fase específica del
crecimiento en el tiempo. Segundo, \textbf{reportar metrias como
Kappa o Accuracy sin FoM infla artificialmente la calidad del modelo}: en un
sistema urbano consolidado donde el 90\% de las celdas permanece sin cambio,
un modelo trivial que predice ``no cambio'' en todas las celdas obtiene
Accuracy de $0.90$ sin capturar ninguna dinámica. Queremos capturar la dinámica, no la estabilidad.

En consecuencia, la validación del WoE-AC, para el presnete trabajo, se realizó en
cinco ventanas quinquenales independientes (2011--2016 a 2015--2020),
reportando FoM, Kappa e IoU simultáneamente. La estabilidad del FoM a través de
las cinco ventanas, no su valor absoluto en una sola, es lo que permite
denotar un modelo que generaliza bien de uno que no lo hace tan bien.

\subsection{Interpretabilidad: el problema de la caja negra}

\textcite{almeida2008stochastic} entrenó una red neuronal feedforward con
backpropagation sobre variables biofísicas e infraestructurales de Piracicaba
(Brasil, 1985--1999); la red aprende pesos sinápticos que alimentan a un CA
estocástico. El modelo reporta concordancia espacial superior a los AC clásicos
de la época, pero los pesos no son traducibles a una pregunta geográfica: no se
puede afirmar que ``la distancia a vialidades contribuyó un $X\%$ a la
transición de esta celda''. La activación de cada neurona es una combinación no
lineal de todas las variables de entrada, sin descomposición legible.

Trabajos posteriores intensifican esta opacidad por dos vías. La primera es la
optimización metaheurística: \textcite{Wang2021} acopla un autómata celular a
un algoritmo genético para generar formas urbanas que maximizan un índice de
vitalidad en Wuhan, donde cada generación produce cromosomas que carecen de
lectura geográfica directa. La segunda es el aprendizaje profundo aplicado a
teledetección y modelado espacial \parencite{Ma2019,Hagenauer2022}, que
requiere GPU y conjuntos de datos masivos para entrenar redes cuyas
representaciones internas no son auditables. La ganancia en exactitud es
real, pero el costo metodológico es alto: el modelo clasifica sin justificar,
y la justificación es exactamente lo que un instrumento de planeación urbana
requiere. Un plan municipal sustentado en ``el modelo lo predijo'' es
jurídica y políticamente más débil que uno sustentado en ``históricamente, la
proximidad al frente urbano contribuye con un Information Value de $X$ a la
probabilidad de transición''.

El método de Pesos de Evidencia provee exactamente esa descomposición. La
elección de WoE en este trabajo no responde a una preferencia estética por la
transparencia: responde al requerimiento R4 de la
Tabla~\ref{tab:requerimientos_modelo_1}, que define la auditabilidad como
condición de uso del modelo, no como propiedad accesoria.

\subsection{Transferibilidad: software propietario y reproducibilidad}

Una fracción significativa de los modelos LUCC más citados depende de software
propietario. Dinamica EGO opera bajo licencia académica restringida; IDRISI y
TerrSet son comerciales; SLEUTH funciona con un código en C disponible pero
cuya instalación en sistemas modernos requiere ajustes no documentados. La
consecuencia práctica es que la reproducibilidad de los resultados publicados
es, en el mejor de los casos, parcial: el lector puede leer el método pero no
ejecutarlo sin acceso a infraestructura específica.

Esta dependencia tiene una consecuencia material poco discutida en la
literatura metodológica: el acceso a licencias y a hardware especializado
introduce una barrera de entrada que condiciona qué grupos académicos pueden
reproducir o extender los modelos publicados. En contextos donde estos recursos
son limitados, la única alternativa práctica es reescribir el pipeline desde
cero. El presente trabajo opta por esta vía: se implementa íntegramente en
Python con dependencias de código abierto (NumPy, SciPy, scikit-learn) sobre
datos públicos (Google Earth Engine, INEGI), satisfaciendo R3 sin
concesiones.

% ============================================================
\section{Pesos de Evidencia y modelos probabilísticos espaciales}
\label{sec:woe-literatura}
% ============================================================

El método de Pesos de Evidencia, desarrollado por \textcite{BonhamCarter1994}
en el contexto de la prospección mineral, se formaliza en el Capítulo~\ref{ch:marco-teorico}
(ecuación~\ref{eq:woe}): para cada clase de una variable espacial discretizada
se calcula un peso bayesiano que cuantifica la asociación entre dicha clase y
la ocurrencia del evento de interés en este trabajo, la transición de suelo
no urbano a urbano. El peso total por celda es la suma de los pesos de las
variables evaluadas en ella, transformada por la función logística para
obtener una probabilidad acotada (ecuación~\ref{eq:woe_sigmoid}); el
Information Value sintetiza el poder discriminatorio de cada variable. La
presente sección no repite la derivación matemática, sino que la sitúa en la
literatura LUCC.

La adaptación al modelado LUCC se consolidó con herramientas como Dinamica EGO,
que combinan mapas de idoneidad WoE con un autómata celular de dos pasos, \textit{expander} y \textit{patcher}, que reproduce la geometría del cambio
observado. La literatura derivada documenta aplicaciones a deforestación
amazónica, expansión agrícola y, en menor medida, crecimiento urbano.

El WoE tiene tres limitaciones estadísticas reconocidas que es necesario
explicitar para evitar atribuirle propiedades que no posee:

\paragraph{Independencia condicional.} La suma de pesos asume que las variables
son condicionalmente independientes dada la ocurrencia del evento. Este
supuesto rara vez se satisface en variables espaciales urbanas, donde
proximidad a vialidades y proximidad al frente urbano están correlacionadas. El
presente trabajo no introduce una corrección formal para esta dependencia;
sigue la práctica estándar de la literatura LUCC, que reporta el sesgo derivado
mediante el contraste empírico con el cambio observado. La validación
multi-período mitiga parcialmente el problema: si el sesgo fuera dominante, el
FoM degradaría sistemáticamente entre ventanas.

\paragraph{Sensibilidad a la discretización.} Los pesos dependen de cómo se
categorice cada variable continua. El presente trabajo adopta una discretización
basada en cuantiles uniformes, que es la convención más reproducible aunque no
necesariamente la óptima en términos de IV.

\paragraph{Multicolinealidad entre drivers.} Variables fuertemente correlacionadas
producen pesos redundantes que inflan la magnitud aparente de la evidencia. Este
trabajo aplica un análisis de componentes principales sobre el conjunto de
variables espaciales antes del cálculo de IV, lo que reduce la redundancia sin
sacrificar interpretabilidad.

% ============================================================
\section{Validación espacial en modelos LUCC}
\label{sec:validacion-literatura}
% ============================================================

La métrica más reportada en LUCC sigue siendo el coeficiente Kappa de Cohen, a
pesar de las críticas explícitas a su uso en mapas de cambio. Kappa mezcla dos
fuentes de error: desacuerdo en cantidad y desacuerdo en asignación, que
tienen significados metodológicos distintos
\parencite{pontius2008comparing}. Esta descomposición separa:
\begin{itemize}
    \item \textbf{Quantity disagreement}: el modelo predice un total de celdas
    de cambio distinto al observado, sin importar dónde.
    \item \textbf{Allocation disagreement}: el modelo predice el total correcto
    de cambio pero lo asigna a celdas incorrectas.
\end{itemize}

Un modelo puede ser perfecto en \textit{quantity} y catastrófico en
\textit{allocation}: predice exactamente el área urbana total pero en lugares
equivocados. El FoM, definido como el índice de Jaccard restringido a las
celdas de cambio, es sensible exactamente a este problema.

La Intersection over Union (IoU) sobre la clase urbana provee información
complementaria: mientras el FoM evalúa la calidad de la predicción del cambio,
la IoU evalúa la calidad de la predicción del estado urbano final. Reportar
ambas evita el sesgo de evaluar solo la dinámica o solo el estado.

El presente trabajo adopta este protocolo de evaluación: FoM como métrica
primaria para el cambio, IoU para el estado, Kappa por compatibilidad con la
literatura previa y un análisis explícito de la descomposición quantity vs.
allocation cuando el FoM es bajo.

% ============================================================
\section{Vacíos específicos en el contexto mexicano}
\label{sec:vacios-mexico}
% ============================================================

La literatura mexicana sobre modelado de crecimiento urbano es escasa pero
reveladora. \textcite{SuarezDelgado2007} construyó el primer estudio de
prospección espacial sistemática para la Zona Metropolitana de la Ciudad de
México: a partir de proyecciones demográficas de CONAPO al 2020, distribuyó
hectáreas adicionales en contornos metropolitanos bajo tres escenarios de
densidad. El trabajo es relevante por su pionería pero no es un modelo CA: no
hay reglas de vecindad, no hay calibración sobre datos históricos de cambio
celular y la validación se hace a nivel de totales de hectáreas, no a nivel de
celda. R1, R4 y la replicabilidad pública no aplican; la única dimensión
satisfecha es R2.

\textcite{RamirezHernandez2021}, en un libro institucional de IIEc-UNAM, aplica
un AC inspirado en SLEUTH a la ZMCM con datos de 1990--2010 y proyecta al
2040. Los parámetros se describen en el texto pero el código no está
disponible y la validación no reporta FoM ni Kappa de forma sistemática. La
obra es valiosa como referencia conceptual pero no como benchmark cuantitativo:
no es posible comparar su desempeño con el rango de
\textcite{pontius2008comparing}.

\textcite{Chihuahua2023} aplica una cadena SVM + regresión logística +
CA-Markov en TerrSet/IDRISI a núcleos urbanos asentados en cinco cuencas
hidrológicas del estado de Chihuahua (Laguna Bustillos y de los Mexicanos,
río Conchos, Ojinaga, río Conchos, Presa de la Colina, río Conchos, Presa el
Granero y río San Pedro), con proyecciones a 2030, 2040 y 2050. La
clasificación SVM sobre mapas de uso de suelo 2010--2020 es reproducible y la
regresión logística aporta interpretabilidad parcial (pseudo $R^2$ de
McFadden superior a $0.2$, ROC aceptable), pero la asignación CA-Markov
ocurre dentro del módulo Land Change Modeler de TerrSet, software
propietario, y la validación cuantitativa se reporta sobre la regresión
logística, no sobre la simulación CA en sí. R4 se cumple parcialmente; R3 no
se cumple.

El cuadrante metodológico que estos tres trabajos definen es claro:

\begin{table}[H]
\centering
\caption{Posicionamiento de la literatura mexicana frente a R1--R4.}
\label{tab:cuadrante-mexico}
\begin{tabular}{lcccc}
\toprule
\textbf{Trabajo} & \textbf{R1} & \textbf{R2} & \textbf{R3} & \textbf{R4} \\
\midrule
\textcite{SuarezDelgado2007}    & ---     & \checkmark & ---        & ---       \\
\textcite{RamirezHernandez2021} & ---     & \checkmark & ---        & parcial   \\
\textcite{Chihuahua2023}        & parcial & \checkmark & ---        & parcial   \\
\midrule
\textbf{Presente trabajo}       & \checkmark & \checkmark & \checkmark & \checkmark \\
\bottomrule
\end{tabular}
\end{table}

Ningún trabajo publicado para una ciudad mexicana cumple simultáneamente con
validación multi-período cuantitativa, reproducibilidad abierta e
interpretabilidad auditable a nivel de celda. Esa es la brecha que el presente
trabajo atiende, y es la razón por la cual la elección de Querétaro como caso
de estudio no es sustituible: aplicar el método sobre una ciudad ya estudiada
por modelos de caja negra no generaría comparación válida con la literatura
existente; aplicarlo sobre una ciudad intermedia sin estudios previos abre el
espacio metodológico.

% ============================================================
\section{Síntesis: posicionamiento del presente trabajo}
\label{sec:sintesis-estado-arte}
% ============================================================

Los cuatro problemas metodológicos analizados (equifinalidad en la
calibración, sobreajuste en la validación, opacidad en la función de
transición y dependencia de software propietario) se cruzan con la geografía
del campo. Las revisiones sistemáticas de \textcite{Sante2010} y
\textcite{Aburas2016} muestran que las aplicaciones documentadas se concentran
en China, Europa y Estados Unidos, mientras que la literatura mexicana
revisada en la sección anterior avanza con modelos heredados o con
herramientas propietarias que limitan la auditoría.

El presente trabajo se ubica en la intersección de tres decisiones:
(i) usar WoE como función de transición, lo que satisface R4 sin renunciar a
la rigurosidad estadística; (ii) validar en cinco ventanas quinquenales
independientes con FoM, Kappa e IoU, lo que satisface R1 y permite distinguir
robustez de sobreajuste; (iii) implementar el pipeline completo en Python con
datos abiertos, lo que satisface R3 y abre la posibilidad de replicación en
otras ciudades intermedias mexicanas. R2, la utilidad para planeación urbana,
se cumple por construcción: los mapas resultantes son raster celda a celda
con probabilidades de transición auditables.

El siguiente capítulo describe el área de estudio y la construcción de la serie
histórica de mapas urbano/no-urbano que alimenta el modelo.
